\documentclass[11pt,a4paper]{article}
\usepackage[utf8]{inputenc}
\usepackage[margin=2.2cm]{geometry}
\usepackage{amsmath,amssymb}
\usepackage{hyperref}
\usepackage{xcolor}
\usepackage{booktabs}
\usepackage{listings}
\usepackage{graphicx}
\usepackage{enumitem}

\hypersetup{
  colorlinks=true,
  linkcolor=blue!50!black,
  urlcolor=blue!50!black,
  citecolor=blue!50!black,
}

\definecolor{codebg}{RGB}{246,248,250}
\definecolor{codecmt}{RGB}{106,115,125}
\definecolor{codekw}{RGB}{215,58,73}
\lstset{
  basicstyle=\ttfamily\footnotesize,
  backgroundcolor=\color{codebg},
  commentstyle=\color{codecmt}\itshape,
  keywordstyle=\color{codekw},
  frame=single,
  framerule=0pt,
  framexleftmargin=4pt,
  framexrightmargin=4pt,
  framextopmargin=2pt,
  framexbottommargin=2pt,
  breaklines=true,
  breakatwhitespace=true,
  showstringspaces=false,
}

\title{FSA v1.5 closed-loop SMC$^2$-MPC: a side-by-side\\
       Python+JAX vs Julia comparison study}
\author{Hand-off report for the next agent}
\date{2026-05-08}

\begin{document}
\maketitle
\tableofcontents

\begin{abstract}
This document is a hand-off for the next agent (or human reader)
who picks up the FSA v1.5 closed-loop comparison work. It explains
what the FSA v1.5 model is, what the \texttt{smc2fc} / \texttt{SMC2FC}
frameworks do, the functional programming paradigm we have moved
to on the Julia side, the comparison problem itself, the tools
developed for it, the bugs found and fixed, where we are today,
and what comes next. All paths are absolute paths inside the
\texttt{python-smc2-filtering-control} repository.
\end{abstract}

\section{Why this report exists}

The FSA model has three Julia codebases and three Python codebases
in this repo (\texttt{version\_1\_*}, \texttt{version\_2\_*},
\texttt{version\_3\_*} on each side). The most relevant pair for
this comparison is:

\begin{itemize}[itemsep=2pt]
  \item \textbf{Python+JAX:} \texttt{version\_1\_5\_Python\_JAX/}
        --- the simplest closed-loop FSA bench, derived from the
        working \texttt{version\_2\_Python\_JAX/} bench, with the
        v1.5 model surface (3-state SDE, 10 estimated parameters,
        3 direct-Gaussian observation channels, no circadian).
  \item \textbf{Julia:} \texttt{version\_1\_5\_Julia/} --- the
        Julia rewrite of the same model surface, written in a
        purely functional style for portability to Lean4 (the
        formal source of truth at \texttt{version\_1\_5\_LEAN/}).
\end{itemize}

The two stacks share the model math (verified bit-equal to 1e-6
through a Julia$\leftrightarrow$Python+JAX$\leftrightarrow$Lean4
121-case differential test). They diverge at the framework layer:
the Python+JAX side uses the mature \texttt{smc2fc} engine; the
Julia side uses its own \texttt{SMC2FC.jl} port plus model-side
GPU kernels in KernelAbstractions.

The user has invested weeks of optimisation in the Python+JAX
side. The Julia side is roughly two days old. This report freezes
the current state of the comparison, the bugs we found and fixed
in this session, and the open work for the next agent.

\section{What is FSA v1.5?}

FSA (Fitness--Stress--Autonomic) is a 3-state stochastic
differential equation model for an athlete's response to
training. State $y = (B, F, A)$:

\begin{itemize}[itemsep=2pt]
  \item $B$ --- aerobic fitness (Banister chronic; bounded
        $[0, 1]$, Jacobi-style state-dependent diffusion).
  \item $F$ --- unified fatigue (Banister acute; $\geq 0$,
        CIR-style diffusion).
  \item $A$ --- autonomic amplitude (Stuart--Landau; $\geq 0$,
        CIR-style diffusion).
\end{itemize}

Drift (\texttt{/day} units):
\begin{align*}
  \mu(B, F)        &= \mu_0 + \mu_B B - \mu_F F - \mu_{FF} F^2 \\
  dB/dt            &= \kappa_B \, (1 + \varepsilon_A A) \, \Phi(t) - B / \tau_B \\
  dF/dt            &= \kappa_F \, \Phi(t) - (1 + \lambda_A A) / \tau_F \, F \\
  dA/dt            &= \mu \cdot A - \eta A^3
\end{align*}

State-dependent diffusion (It\^o):
\begin{align*}
  \sigma_B(B) \, dW_B &= \sigma_B \sqrt{B(1-B)} \, dW_B \\
  \sigma_F(F) \, dW_F &= \sigma_F \sqrt{F} \, dW_F \\
  \sigma_A(A) \, dW_A &= \sigma_A \sqrt{A} \, dW_A
\end{align*}

The single exogenous control is $\Phi(t)$, the training-strain
rate ($\geq 0$). The closed-loop MPC's job is to pick $\Phi$ to
maximise the time-integral of $A$ subject to a soft fatigue
constraint ($\max(F - F_{\max}, 0)^2$ barrier).

\subsection{What makes v1.5 special}

\begin{itemize}[itemsep=2pt]
  \item \textbf{14 dynamics parameters} in the v1.5 basis;
        $\kappa_B$ replaced by $B_\infty = \kappa_B \tau_B$ and
        $\kappa_F$ replaced by $F_\infty = \kappa_F \tau_F$ for
        better identifiability. The basis-rotation adapter is in
        \texttt{simulation.params\_v15\_to\_v1\_nt} (Julia) /
        \texttt{simulation.params\_v15\_to\_v1} (Python+JAX).
  \item \textbf{10 estimated parameters} (the filter's decision
        variables); 4 dynamics parameters pinned at truth
        ($\tau_B$, $\eta$, $\varepsilon_A$, $\mu_{FF}$) per the
        FIM-rank decision; 3 obs-noise scales pinned at 0.005.
  \item \textbf{3 direct-Gaussian observation channels}: $y^{obs}_B
        = B + \mathcal{N}(0, \sigma_{B,obs}^2)$ and same for $F$,
        $A$. No HR / sleep / stress / steps channels (that's the
        v2 surface). No circadian forcing.
  \item Time grid: 60 minutes per bin (24 bins per day) by default.
        v2 uses 15 minutes per bin (96 bins per day).
\end{itemize}

\section{What \texttt{smc2fc} / \texttt{SMC2FC} does}

Both stacks implement the same statistical pipeline:

\begin{enumerate}[itemsep=2pt]
  \item \textbf{Outer SMC$^2$ filter}: tempered sequential Monte Carlo
        over the 10-dim unconstrained parameter cloud $\theta$.
        At each tempering level $\lambda \in [0, 1]$ the targets
        graduate from prior to posterior. Resample
        (systematic) $\to$ Hamiltonian Monte Carlo (HMC) moves
        with finite-difference gradients $\to$ next $\lambda$.
        Returns a posterior cloud over $\theta$ for each rolling
        filter window.
  \item \textbf{Inner particle filter (PF)}: per $\theta$ candidate,
        an inner PF computes $\log p(y_{1:T} \mid \theta)$ by
        propagating a state-particle cloud forward through the SDE
        and reweighting by the observation likelihood. This is the
        marginal log-likelihood the outer SMC$^2$ uses.
  \item \textbf{Forward-simulation controller}: NOT a particle
        filter. The controller takes the filter-derived posterior
        mean parameters $\hat\theta$ and the filter-derived smoothed
        end-of-window state $\hat y$. It then runs an SMC$^2$ over
        a low-dimensional schedule parameterisation (8 RBF
        coefficients $\theta_{ctrl}$): for each candidate
        $\theta_{ctrl}$, decode the 8 anchors into a
        per-bin $\Phi$ schedule, simulate the SDE forward
        $n_{\text{inner}}$ times under common-random-numbers (CRN)
        noise, accumulate the cost
        $J = -\!\int A \, dt + \lambda_F \!\int (F - F_{\max})_+^2 \, dt$,
        and use the negative cost as a tempering target. Output is
        the posterior-mean $\theta_{ctrl}^\star$, decoded to a
        $\Phi$-schedule for the plant's next stride.
\end{enumerate}

\textbf{Critical: the controller does not run a particle filter
internally. It runs deterministic forward simulations
($n_{\text{inner}}$ noise trials per $\theta_{ctrl}$) using the
filter's $\hat\theta$ and $\hat y$.} This was an explicit design
question raised during the comparison and the answer is the same
on both stacks: both controllers correctly forward-simulate.

\section{Functional programming paradigm in Julia}

\subsection{Why we moved to functional}

The \texttt{version\_3\_julia/} stack was a Python-port-style Julia
implementation: mutable structs, \texttt{!}-suffix mutating methods,
threaded \texttt{MersenneTwister} state, history dicts pushed onto
across stride boundaries. It accumulated subtle bugs that were hard
to trace, and the user explicitly rejected it (``too much for a
human to understand'').

\texttt{version\_1\_5\_Julia/} is a clean re-design. The five rules:

\begin{enumerate}[itemsep=2pt]
  \item \textbf{No mutable structs in the public API.} v1.5 has
        exactly one struct, \texttt{PlantState}, with two fields:
        an \texttt{SVector\{3, Float64\}} and an \texttt{Int}. State
        updates produce a NEW \texttt{PlantState}, not a mutation
        of the old one.
  \item \textbf{No \texttt{!}-suffix mutating methods.} Every public
        function returns new values. \texttt{plant\_step(state,
        Phi\_t, params, dt, key)} returns
        \texttt{(state=new\_state, obs=...)}; the caller can build
        history by collecting returned values.
  \item \textbf{Explicit RNG keys, JAX-style.} Pass an explicit
        \texttt{key::UInt64} to every randomised function. Internal
        sub-keys are derived deterministically via
        \texttt{hash((key, :sub\_purpose))}. Two callers with the
        same \texttt{(state, key)} get the same answer, period.
  \item \textbf{Pure facade over the inherently imperative GPU
        kernel.} KernelAbstractions kernels mutate preallocated
        \texttt{CuArrays} by design. v1.5 wraps every kernel-touching
        public function in a pure facade: \texttt{gpu\_log\_density(
        target, U, grid\_obs, key) $\to$ Vector\{Float64\}} allocates
        a fresh output, dispatches the kernel, returns. Mutation is
        encapsulated; callers never see the \texttt{CuArray}.
  \item \textbf{I/O isolated to the bench's top-level
        \texttt{main()}.} JLD2 + PNG writes happen exactly once at
        the end of the bench. Logging at the bench level, not inside
        \texttt{plant\_step} / \texttt{propagate} /
        \texttt{obs\_log\_weight}.
\end{enumerate}

\subsection{The \texttt{SMC2FC\_functional} library --- now the default}

\textbf{Status update (2026-05-08):} the functional rewrite at
\texttt{julia/SMC2FC\_functional/} is now \emph{the default Julia
framework}. The original port has been renamed to
\texttt{julia/SMC2FC\_DEPRECIATED\_USE\_FUNCTIONAL/} (kept as a
back-compat reference for the three-library comparison harness only).
A bit-identical migration test is documented in
\S\ref{sec:bug_migration} and recorded in
\texttt{julia/SMC2FC\_functional/docs/MIGRATION\_RESULT.md}.

The functional library has the same algorithmic contract as the
deprecated port, but:

\begin{itemize}[itemsep=2pt]
  \item Top-level functions take immutable inputs and return
        immutable outputs.
  \item Pre-allocated buffers live behind an immutable
        \texttt{Workspace} container.
  \item No \texttt{!}-mutating functions exposed publicly.
  \item Multiple-dispatch-driven specialisation; idiomatic Julian.
  \item Google-style docstrings on every file and every function.
\end{itemize}

The drop-in bench at\\
\texttt{julia/SMC2FC\_functional/benchmarks/gpu\_drop\_in/bench\_dropin.jl}\\
is byte-equivalent to the v1.5 GPU bench except for one import line
(\texttt{using SMC2FC\_functional: run\_tempered\_smc\_gpu}). The
three-library cross-check harness lives at\\
\texttt{julia/SMC2FC\_functional/benchmarks/compare\_three\_libraries.jl}
and the seven verification gates are written up in
\texttt{julia/SMC2FC\_functional/docs/GATE\_RESULTS.md}.

\section{The comparison problem}

The user's question: \emph{do the Python+JAX and Julia v1.5 stacks
produce qualitatively the same closed-loop MPC behaviour, given
they implement the same model under the same statistical
framework?}

This is harder than it sounds for several reasons:

\begin{enumerate}[itemsep=2pt]
  \item Different RNG implementations (JAX \texttt{PRNGKey} vs
        Julia \texttt{StableRNG} vs hash-derived sub-keys) mean
        the plant trajectories are NOT bit-identical even at the
        same seed. Comparison must be statistical: mean $\int A$,
        F-violation rate, posterior MSE to truth, schedule shape.
  \item Different inner-PF algorithms: Python+JAX uses
        \texttt{smc2fc.filtering.gk\_dpf\_v3\_lite} (a Gaussian
        kernel, OT-regularised, differentiable PF); Julia uses a
        custom segmented-PF KernelAbstractions kernel. Both are
        valid, but they have different particle-degeneracy
        regimes.
  \item Different default particle counts: Python+JAX bench
        defaulted to $N_{\text{smc}}=128 / N_{\text{inner}}=32$ on
        the controller; the Julia bench defaulted to
        $N_{\text{smc}}=256 / N_{\text{inner}}=64$. We brought
        these into parity.
  \item No prior baseline: no constant-$\Phi=1$ baseline trajectory
        was being plotted alongside the MPC, so ``does the
        controller actually beat doing nothing?'' was hard to
        answer at a glance.
  \item Different output formats: Python writes \texttt{.npz}, Julia
        writes \texttt{.jld2} (HDF5).
\end{enumerate}

\subsection{Tools developed for the comparison}

\textbf{1. JAX GPU profiler}
(\texttt{src/profile\_gpu\_python.py}). Direct port of the existing
v2 Julia profiler. Times the GK-DPF v3-lite \emph{filter kernel}
and the controller cost \emph{kernel} in isolation, with
\texttt{jax.block\_until\_ready} as the synchronisation point. Reports
median time per call, threads-in-flight, approximate FLOPs,
effective TFLOPS, percentage of the RTX-5090's 104.8 TFLOPS fp32
peak, GPU memory before and after, and host-sync count.

\textbf{2. v1.5 Julia GPU profiler}
(\texttt{src/profile\_gpu\_julia\_v15.jl}). Same surface, same
metrics, adapted to v1.5's GPU types (\texttt{FSAGPUTarget},
\texttt{FSAv1ControlGPUTarget}).

\textbf{3. Per-stride telemetry CSV} (added to both bench
scripts). Each stride writes one row to \texttt{per\_stride.csv}
with: \texttt{stride}, \texttt{t\_wall\_s},
\texttt{n\_temp\_filter}, \texttt{n\_temp\_ctrl},
\texttt{daily\_phi}, \texttt{A\_mean\_so\_far}, \texttt{B\_end},
\texttt{F\_end}, \texttt{A\_end}.

\textbf{4. Paired-run launcher}
(\texttt{src/run\_v15\_T28d\_compare.sh}). Single shell script
that:
\begin{enumerate}[itemsep=2pt]
  \item Runs both microbench profilers (writes
        \texttt{julia\_profile.log} + \texttt{python\_profile.log}).
  \item Starts \texttt{nvidia-smi} sampling at 1 Hz to a CSV.
  \item Runs the Julia closed-loop bench with matched config.
  \item Stops + restarts \texttt{nvidia-smi} for the Python run.
  \item Runs the Python+JAX closed-loop bench, same matched config.
\end{enumerate}
Output lands under
\texttt{outputs/v15\_julia\_vs\_python\_T<T>d\_seed<S>/\{julia,python\}/}.

\textbf{5. Three-way comparison script}
(\texttt{src/compare\_v15\_julia\_vs\_python.py}). Loads both
bench outputs, both \texttt{nvidia-smi} CSVs, both profile logs.
Produces:
\begin{itemize}[itemsep=2pt]
  \item \texttt{summary.txt} --- the diagnostic table (baseline as
        gate, then macro \texttt{nvidia-smi} telemetry, then micro
        profiler numbers).
  \item \texttt{comparison.png} --- 6-panel three-way diagnostic
        plot. Python MPC blue, Julia MPC red, baseline grey
        dashed. Panels: $A$ trajectory, applied daily $\Phi$,
        posterior medians for two parameters with truth lines,
        GPU SM utilisation over wall time, per-stride wall time
        bars.
\end{itemize}
JLD2 files are read with \texttt{h5py} and the 2D and 3D arrays are
transposed (Julia is column-major; HDF5 readers in row-major order
swap dimensions).

\section{Bugs found and fixed}

We worked top-down through the comparison and found four
substantive issues, two on Python+JAX and two on Julia.

\subsection{Python: shrinking planning horizon (A1)}

\textbf{File}: \texttt{version\_1\_5\_Python\_JAX/tools/bench\_smc\_full\_mpc\_fsa\_v15.py},
old line 287.

The replan block computed
$t_{\text{remaining\_days}} = (n_{\text{strides}} - s) \cdot
\text{STRIDE\_BINS} / \text{BINS\_PER\_DAY}$ and passed
\texttt{T\_total=t\_remaining\_days} into
\texttt{build\_control\_spec()}. As the bench progresses, the
horizon shrinks; at stride 14 of T=28d the controller sees only
~14 days ahead. v2's bench instead always plans the FULL
\texttt{T\_total\_days} at every replan
(\texttt{bench\_smc\_full\_mpc\_fsa.py:386-388}). v1.5 Julia uses
\texttt{H\_plan\_bins = T\_total\_bins} (also full horizon).

\textbf{Symptom}: the controller could not justify a "build $B$
then sprint" Banister strategy because the time-to-horizon was
constantly shrinking.

\textbf{Fix}: pass \texttt{T\_total=float(args.T\_days)} (constant)
at every replan.

\subsection{Python: mid-array splice instead of overwrite (A2)}

\textbf{File}: same, old lines 317--320.

The new plan was spliced INTO THE MIDDLE of the existing
\texttt{daily\_phi\_plan} array using an absolute-current-day
offset. Combined with the shrinking horizon (A1), this produced a
sawtooth $\Phi$ pattern: the front of the plan got partially
overwritten on each replan, the back stayed stale.

v2 overwrites the whole plan: \texttt{daily\_phi\_plan =
sched\_per\_day.astype(np.float64)} +
\texttt{last\_replan\_stride = s + 1}.

\textbf{Fix}: same overwrite-and-reset pattern in v1.5 Python.

\subsection{Python: filter degeneracy from bootstrap with GK-DPF (B)}

\textbf{File}: \texttt{version\_1\_5\_Python\_JAX/models/fsa\_high\_res/estimation.py},
function \texttt{\_propagate\_fn\_framework}.

This was the structural bug. The \texttt{smc2fc.filtering.gk\_dpf\_v3\_lite}
filter is designed for \emph{guided} (Kalman-fused) proposals: the
\texttt{propagate\_fn} contract expects to return both the new state
\emph{and} a Radon--Nikodym correction \texttt{pred\_lw} that the
framework adds to the per-bin weight increment.

v1.5 Python+JAX's \texttt{\_propagate\_fn\_framework} returned
\texttt{pred\_lw = 0.0} (the bootstrap convention: proposal IS
dynamics, no correction needed). Mathematically valid but
fundamentally incompatible with GK-DPF v3-lite's
weight-smoothing kernel + OT regularisation: the framework's
machinery pushed particles toward biased fixed points. Smoking
gun: at higher particle counts $N=1024 / K=800$ the bootstrap
variant got WORSE (posterior MSE 1.41 $\to$ 31.16), the classic
``structural mis-spec'' signature.

\textbf{Fix}: Kalman fusion. Mirrors v2's
\texttt{propagate\_fn} pattern exactly. v1.5's obs model is the
simplest possible Kalman case ($H = I$, $b = 0$,
$R = \text{diag}(\sigma_{*,obs}^2)$, all channels always
present), so the sequential-scalar update reduces to per-channel
1D Kalman:

\begin{align*}
  \mu_{\text{prior}} &= y + dt \cdot \text{drift} \\
  P_{\text{prior}}   &= \text{diag}(\sigma_B^2 \, B(1{-}B) \, dt,\
                                       \sigma_F^2 \, F \, dt,\
                                       \sigma_A^2 \, A \, dt) \\
  \text{innov}_i     &= y_{\text{obs},i} - \mu_i \\
  S_i                &= P_{ii} + \sigma_{i,obs}^2 \\
  K_i                &= P[:, i] / S_i \\
  \mu                &\mathrel{+}= K_i \cdot \text{innov}_i \\
  P                  &\mathrel{-}= K_i \otimes P[i, :] \\
  \log p_i           &= -\tfrac12 \log(2\pi S_i) - \tfrac{\text{innov}_i^2}{2 S_i}
\end{align*}

Sample $x_{\text{new}} \sim \mathcal{N}(\mu_{\text{fused}},
P_{\text{fused}})$ via Cholesky. Return \texttt{pred\_lw =
log\_pred\_total - obs\_ll(y\_new)} so the framework's
\texttt{obs\_log\_weight\_fn} re-add cancels properly.

\textbf{Result} (T=7d, seed=42, $N=32$, $K=200$):

\begin{center}
\begin{tabular}{lrrr}
  \toprule
                     & Before  & After    & Improvement \\
  \midrule
  posterior MSE      & 1.41    & 0.0188   & $\sim$74$\times$ \\
  \bottomrule
\end{tabular}
\end{center}

The 10 estimated parameters' posterior medians now hug truth
tightly (was: $\tau_F$ stuck at 3 vs truth 7; sigmas drifting
upward unbounded).

\subsection{Python: \texttt{init\_state} cheat (C)}

\textbf{File}: same bench file, around old line 294.

The bench was passing \texttt{init\_state\_dict =
dict(B=plant\_state.bfa[0], ...)} directly to the controller. This
is a closed-loop violation: the controller should know only what
the filter knows (its posterior estimate of the state), not the
plant's true state.

\textbf{Fix}: extract a smoothed end-of-window state from the
filter posterior cloud via \texttt{log\_density\_factory.extract\_state\_at\_step},
mirroring v2's pattern (line 366--367). Use that as
\texttt{init\_state}.

\subsection{Julia: per-tempering-level allocation churn (perf)}

\textbf{File}: \texttt{version\_1\_5\_Julia/models/fsa\_high\_res/gpu\_pf.jl},
function \texttt{parallel\_hmc\_one\_move}.

The functional HMC variant allocates roughly ten $(M, d)$
matrices per call. With 13 strides $\times$ 5 tempering levels
$\times$ 3 HMC moves per level the bench does $\sim$200 HMC
calls, allocating a few MB total --- not catastrophic, but a
real per-call constant.

\textbf{Fix}: added a mutating
\texttt{parallel\_hmc\_one\_move!} variant that reuses two
persistent $(M, d)$ buffers across the whole leapfrog. The bench
imports both; \texttt{run\_outer\_smc} now calls the mutating
variant.

\textbf{Result} (T=2d open-loop smoke, $N=16$, $K=100$): wall
time 96.9~s $\to$ 92.0~s ($\sim$5\% faster at this small N). All
six GPU PF smoke tests still pass.

\subsection{Julia: optional Liu--West $\theta$-cloud shrinkage}

\textbf{File}: same plus
\texttt{version\_1\_5\_Julia/tools/bench\_smc\_full\_mpc\_fsa\_gpu.jl}.

Added a CLI flag \texttt{--liu-west-a} (default 0.0 = disabled).
When set in $(0, 1)$ --- e.g.\ 0.97 --- applies the standard
Liu--West shrinkage between resample and HMC inside
\texttt{run\_outer\_smc}:

\[
  \theta_i \;\leftarrow\; a \cdot \theta_i \;+\; (1-a) \cdot \overline\theta
                  \;+\; \sqrt{1 - a^2} \cdot \text{column\_std} \cdot \xi,
                  \qquad \xi \sim \mathcal{N}(0, I)
\]

Helps maintain $\theta$-cloud diversity at low $N$, where pure
tempered HMC can collapse. Off by default so existing
reproducibility is preserved.

\subsection{Julia: per-segment GPU\,$\to$\,CPU transfers in
            \texttt{gpu\_log\_density} (perf, big)}\label{sec:bug_segment}

\textbf{File}: \texttt{version\_1\_5\_Julia/models/fsa\_high\_res/gpu\_pf.jl},
function \texttt{gpu\_log\_density}, segment loop.

The segment loop ran a per-chain log-likelihood accumulator
\emph{on the host}:

\begin{lstlisting}
for seg in 1:n_segments
    target.propagate_kernel(...)
    KernelAbstractions.synchronize(CUDABackend())
    @match (seg, n_segments) begin ... end
    # ── HOT-LOOP HOST READ (anti-pattern) ────
    log_max_cpu = Array(view(bufs.log_max, 1:M))   # GPU -> CPU
    log_z_cpu   = Array(view(bufs.log_z,   1:M))   # GPU -> CPU
    @inbounds for m in 1:M
        log_lik_acc[m] += Float64(log_max_cpu[m] + log_z_cpu[m] - log_K)
    end
end
\end{lstlisting}

Per call: 12 GPU\,$\to$\,CPU transfers (M floats $\times$ 2 fields
$\times$ 6 segments) + 6 redundant
\texttt{KernelAbstractions.synchronize()} calls. Each transfer
forces an implicit stream sync that stalls the GPU between segments.

\textbf{Fix}: a Float64 GPU-resident accumulator
\texttt{log\_lik\_acc\_gpu::CuArray\{Float64,1\}} added to
\texttt{FSAGPUTarget}, plus a tiny
\texttt{gpu\_log\_lik\_accumulate\_kernel!} that writes
\texttt{log\_lik\_acc[m] += Float64(log\_max[m] + log\_z[m] - log\_K)}
on the GPU. The end-of-function \texttt{Array(view(...))} syncs the
stream once and returns the M-vector to the host. Removed the
redundant per-segment syncs at the same time.

This change alone was a small win at the tested config ($\sim 3$\,ms
out of 81\,ms). The headline cost was found to live elsewhere ---
see \S\ref{sec:bug_ot}.

\subsection{Julia: OT-rescue per-chain Julia loop (perf, headline)}\label{sec:bug_ot}

\textbf{File}: \texttt{julia/SMC2FC/src/Filtering/GPUSegmentedPF.jl}
(framework, now \texttt{SMC2FC\_DEPRECIATED\_USE\_FUNCTIONAL}; same
code copied verbatim into \texttt{SMC2FC\_functional}), function
\texttt{gpu\_ot\_blend\_chain!}.

The OT (optimal-transport) sigmoid-blend rescue inside
\texttt{run\_segmented\_smc\_step!} wraps the per-chain Sinkhorn
in a Julia for-loop over $M$ chains, which inside its body does

\begin{lstlisting}
log_w_arr = Array(chain_log_w)    # GPU -> CPU memcpy + sync (per chain)
... softmax on CPU ...
b_gpu     = CuArray(b_cpu)        # CPU -> GPU memcpy + sync (per chain)
\end{lstlisting}

At $M=32$ chains $\times$ 5 segments per \texttt{gpu\_log\_density}
call, that is $\sim$320 PCIe round-trips \emph{per call}, every one
of them stalling the GPU. The actual GPU compute under the rescue
(Sinkhorn iterations + barycentric projection) is sub-millisecond.

\textbf{A/B microbench} (RTX 5090, $N_{\text{smc}}=32$,
$K_{\text{pf}}=200$, $T_{\text{steps}}=24$):

\begin{center}
\begin{tabular}{lr}
  \toprule
  \texttt{ot\_max\_weight = 0.01} (was default) & 84.78\,ms / call \\
  \texttt{ot\_max\_weight = 0.0}  (now default) &  0.73\,ms / call \\
  \midrule
  speed-up                                       & $116\times$ \\
  \bottomrule
\end{tabular}
\end{center}

\textbf{Closed-loop A/B} (T=2d, seed=42, $N=16$, $K=100$):

\begin{center}
\begin{tabular}{lrr}
  \toprule
  metric                                   & OT on    & OT off (now default) \\
  \midrule
  total wall time                          & 159.5\,s & 12.5\,s              \\
  mean $A$                                 & 0.0928   & 0.0928               \\
  final $\hat x$                           & (0.059, 0.273, 0.088)
                                                       & (0.059, 0.273, 0.089) \\
  posterior medians (10 params, identical) & 7 of 10  & 7 of 10              \\
  posterior medians (3 params)             &
       \multicolumn{2}{c}{differ $<22\%$ rel., inside the IQR band} \\
  \bottomrule
\end{tabular}
\end{center}

\textbf{Fix}: added \texttt{--ot-max-weight} CLI flag to the bench
(default 0.0, OT off), and flipped the \texttt{FSAGPUTarget}
constructor's \texttt{ot\_max\_weight} keyword default from 0.01
to 0.0. Heavy warning comments at both spots explaining the
slowdown, the microbench and closed-loop numbers, when to turn
it back on (low-ESS / hard obs models), and the proper fix
(a future framework rewrite of \texttt{gpu\_ot\_blend\_chain!}
as a single batched-across-chains GPU kernel).

\textbf{Why this is safe at v1.5}: the writeup's earlier section
on Python's bug (B) showed the GK-DPF v3-lite framework expects
guided proposals and gets them via Kalman fusion --- the OT rescue
is then unnecessary on the strong direct-Gaussian 3-channel obs
model. Python+JAX runs without an explicit OT rescue and matches
truth fine. Julia at v1.5 is in the same regime.

\subsection{Julia framework migration: bit-identical}\label{sec:bug_migration}

The migration of the v1.5 GPU bench from the old \texttt{SMC2FC}
framework to \texttt{SMC2FC\_functional} is exactly one line on the
bench side:

\begin{lstlisting}
- using SMC2FC: run_tempered_smc_gpu
+ using SMC2FC_functional: run_tempered_smc_gpu
\end{lstlisting}

The drop-in artefact at\\
\texttt{julia/SMC2FC\_functional/benchmarks/gpu\_drop\_in/bench\_dropin.jl}\\
is byte-identical to the v1.5 GPU bench except for that import line
and the \texttt{REPO\_ROOT} pointer. \texttt{SMC2FC\_functional}'s
\texttt{Control/GPUControlSMC.jl} was a verbatim COPY+DOC of the
deprecated framework's --- only the file-header block-comment was
converted to a Julia docstring. So at matched config and seed, the
two frameworks must produce the same output, modulo JIT-warmup
wall-time variance.

Head-to-head test at T=2d, seed=42, $N=16$, $K=100$:

\begin{center}
\begin{tabular}{lrrr}
  \toprule
  quantity                            & deprecated \texttt{SMC2FC}
                                      & \texttt{SMC2FC\_functional}
                                      & max\,$|$diff$|$ \\
  \midrule
  wall time                           & 12.5\,s  & 12.9\,s  & +0.4\,s (JIT) \\
  trajectory $B / F / A$ (all bins)   & ---      & ---      & 0.000e+00     \\
  applied $\Phi$ per bin              & ---      & ---      & 0.000e+00     \\
  baseline trajectory                 & ---      & ---      & 0.000e+00     \\
  all 10 posterior medians            & identical& identical& 0.000e+00 each\\
  posterior particle clouds           & ---      & ---      & bit-identical \\
  \bottomrule
\end{tabular}
\end{center}

Every numeric in the JLD2 file is byte-identical.
\texttt{SMC2FC\_functional} is the default framework going forward.

\section{Where we are now}

\subsection{State of play before today}

The pre-fix headline at T=7d, seed=42, matched config
$N_{\text{smc}}=32$, $K_{\text{pf}}=200$ was:

\begin{center}
\begin{tabular}{lrrr}
  \toprule
                                        & Baseline ($\Phi=1$) & Python+JAX MPC & Julia MPC \\
  \midrule
  mean $A$ (last 7 days)                & 0.0935              & 0.0808         & 0.0815    \\
  improvement vs baseline               & ---                 & $-13.6\%$      & $-12.8\%$ \\
  F-violation rate                      & 0.0                 & 0.0            & 0.0       \\
  posterior MSE to truth (final window) & ---                 & 0.0188         & 0.0868    \\
  total wall time (min)                 & ---                 & 5.3            & 37.2      \\
  mean GPU utilisation (\%)             & ---                 & 40.0           & 24.8      \\
  filter kernel time / call (ms)        & ---                 & 4.7            & \textbf{81.0} \\
  controller cost / call (ms)           & ---                 & 23.4           & 0.5       \\
  \bottomrule
\end{tabular}
\end{center}

\subsection{State of play after today's perf work}

After the \texttt{gpu\_pf.jl} GPU-accumulator fix
(\S\ref{sec:bug_segment}), the OT-rescue default flip to off
(\S\ref{sec:bug_ot}), and the framework migration to
\texttt{SMC2FC\_functional} (\S\ref{sec:bug_migration}), the Julia
microbench at the same matched config now reports:

\begin{center}
\begin{tabular}{lrrl}
  \toprule
  metric (RTX 5090, $N_{\text{smc}}=32$, $K_{\text{pf}}=200$,
                     $T_{\text{steps}}=24$)
                                       & Python+JAX & Julia        & ratio \\
  \midrule
  filter kernel time / call (ms)       & 4.7        & \textbf{0.9} & \textbf{Julia $\sim 5\times$ faster} \\
  controller cost / call (ms)          & 23.4       & 0.4          & Julia $\sim 60\times$ faster \\
  \bottomrule
\end{tabular}
\end{center}

i.e.\ the microbench has flipped from \emph{Julia 17$\times$ slower}
(81\,ms) to \emph{Julia 5$\times$ faster} (0.9\,ms) on the filter
kernel, simply by removing the per-segment GPU\,$\to$\,CPU host
read pattern and disabling the OT rescue's per-chain Julia loop.

\textbf{The new headline numbers are nuanced (read carefully):}

\begin{itemize}[itemsep=2pt]
  \item \textbf{Both filters remain algorithmically sound.} The
        Python+JAX MSE 0.0188 vs Julia 0.0868 difference is the
        same maturity gap as before --- not affected by today's
        perf work, which is purely about wall time.
  \item \textbf{The neither-beats-baseline-at-T=7d finding is
        unchanged.} The H3 controller-cost-design issue is
        unaffected by perf. The longer-horizon test (T=14d / T=28d)
        is still the right next experiment.
  \item \textbf{Julia is now \emph{faster per call} than Python+JAX
        --- but the closed-loop wall-time race has moved.} The
        81\,ms filter cost is gone, but the \emph{full T=7d
        closed-loop bench numbers above were captured with the
        slow Julia path and have not been re-measured}. We expect
        a $\sim$10$\times$ closed-loop speed-up on Julia at
        matched config, putting Julia close to or under Python+JAX's
        5.3\,min; that is the next thing to verify.
  \item \textbf{Caveat: matched config is not maxed config.} The
        $N_{\text{smc}}=32 / K_{\text{pf}}=200$ config above is
        what Python+JAX was tuned for. Julia, with its
        $\sim 5\times$ faster per-call cost, has substantial
        headroom on the RTX 5090 --- both compute and memory.
        Bumping particle counts to a Julia-saturation tier (and
        re-running T=14d / T=28d there) is the next-cheapest
        scientific test, and is the focus of the next plan.
        See \S\ref{sec:next}.
\end{itemize}

\section{What comes next}\label{sec:next}

In rough priority order:

\paragraph{1. Migrate to \texttt{SMC2FC\_functional}.\hspace{0.2em}\textsc{[done]}}
\texttt{julia/SMC2FC\_functional/} is now the default Julia
framework. The deprecated \texttt{julia/SMC2FC/} has been renamed
to \texttt{julia/SMC2FC\_DEPRECIATED\_USE\_FUNCTIONAL/} and is
retained only for the three-library comparison harness. The
migration is bit-identical; see \S\ref{sec:bug_migration} and
\texttt{julia/SMC2FC\_functional/docs/MIGRATION\_RESULT.md}.

\paragraph{2. Optimise Julia GPU utilisation.\hspace{0.2em}\textsc{[partially done; new plan filed]}}

\textbf{Done in this session:}
\begin{itemize}[itemsep=2pt]
  \item Per-segment GPU\,$\to$\,CPU host read removed from
        \texttt{gpu\_log\_density}'s segment loop, replaced by a
        GPU-resident accumulator + a tiny accumulation kernel
        (\S\ref{sec:bug_segment}). Removed the redundant
        per-segment \texttt{KernelAbstractions.synchronize()}
        calls that the old code needed.
  \item OT-rescue per-chain Julia loop turned off by default
        (\S\ref{sec:bug_ot}). The headline filter kernel cost
        dropped from 81\,ms / call to 0.9\,ms / call.
  \item Mutating HMC variant kept (the earlier session's small
        win at low N).
\end{itemize}

\textbf{Open:} the closed-loop wall time at T=7d / T=14d / T=28d
has not been re-measured with the new fast path, and Julia is
\emph{not} yet running at a config that saturates the RTX 5090's
104.8 TFLOPS fp32 peak / 32\,GB memory. The new plan filed at
\texttt{claude\_plans/Julia\_T28d\_max\_5090\_saturation\_2026-05-08.md}
covers:

\begin{itemize}[itemsep=2pt]
  \item Re-running the headline microbench at $N_{\text{smc}}=256$,
        $K_{\text{pf}}=400$ (v2 saturation tier) and at
        $N_{\text{smc}}=512$ / $K=800$ to find the GPU's compute
        and memory ceiling.
  \item \texttt{nsys profile} on a real bench run to attribute
        the remaining microseconds to specific kernels.
  \item Picking the largest config that fits in memory at T=28d
        and re-running the full closed-loop comparison at that
        scale, then re-running the section-9 comparison harness
        end-to-end.
  \item Re-introducing OT rescue \emph{only} as a future
        framework rewrite of \texttt{gpu\_ot\_blend\_chain!} into
        a single batched-across-chains GPU kernel (out of scope
        for this plan).
\end{itemize}

\paragraph{3. Run T=14d / T=28d to test H3.}
The user's reference v2 plots ramp $\Phi$ UP at longer horizons
(T=14d: mean $A$ 0.122 vs baseline 0.081, $+50\%$; T=28d: 0.157
vs 0.098, $+60\%$). The hypothesis is that v1.5's controller will
similarly find the ramp-up optimum at horizons long enough for
the ``build $B$ then sprint'' payoff to dominate. Now that the
filter kernel is $\sim$90$\times$ faster, T=28d at the saturated
config is wall-time-feasible (estimated under 30\,min on Julia).
This experiment is wrapped into the plan at the end of item 2.

\paragraph{4. Cost-function variants.}
If T=14d / T=28d still picks rest, try:
\begin{itemize}[itemsep=2pt]
  \item Re-introduce $\lambda_\Phi \!\int \Phi^2 \, dt$ penalty
        (currently $\lambda_\Phi = 0$ on both stacks by default
        --- the cost expression supports it).
  \item Add a soft ``minimum mean-$A$ over planning window''
        floor.
  \item Investigate whether the v2 G1 reparametrisation around
        $(A_{TYP}, F_{TYP})$ shifts the bifurcation balance in
        a way that makes ramp-up the dominant optimum at all
        horizons.
\end{itemize}

\paragraph{5. Multi-seed robustness.}
All comparison runs to date are single-seed. If a T=28d run flips
H3 cleanly, replicate at $\geq 3$ seeds before declaring victory.

\section{How to run the comparison end-to-end}

The full pipeline is automated by the launcher:

\begin{lstlisting}
conda activate comfyenv
bash /home/ajay/Repos/python-smc2-filtering-control/version_1_5_Python_JAX/tools/launchers/run_v15_T28d_compare.sh 7 42
# args: T_DAYS, SEED. Defaults: T=28, seed=42.
\end{lstlisting}

This produces (under the \texttt{outputs/v15\_julia\_vs\_python\_T<T>d\_seed<S>/}
parent dir):

\begin{itemize}[itemsep=2pt]
  \item \texttt{julia\_profile.log}, \texttt{python\_profile.log}
  \item \texttt{julia/}: \texttt{data.jld2},
        \texttt{manifest.json}, \texttt{per\_stride.csv},
        \texttt{gpu\_telemetry.csv}, two PNGs (state traces, param
        traces), \texttt{bench.log}.
  \item \texttt{python/}: \texttt{trajectory.npz},
        \texttt{manifest.json}, \texttt{per\_stride.csv},
        \texttt{gpu\_telemetry.csv}, two PNGs, \texttt{bench.log}.
\end{itemize}

Then run the comparison script:

\begin{lstlisting}
cd /home/ajay/Repos/python-smc2-filtering-control/version_1_5_Python_JAX
JAX_ENABLE_X64=True PYTHONPATH=.:.. \
  python tools/compare_v15_julia_vs_python.py \
  /home/ajay/Repos/python-smc2-filtering-control/outputs/v15_julia_vs_python_T7d_seed42
\end{lstlisting}

This writes \texttt{summary.txt} and \texttt{comparison.png} into
the parent dir.

\section{How to profile each stack in isolation}

\begin{lstlisting}
# Julia microbench (filter kernel + controller cost kernel)
cd /home/ajay/Repos/python-smc2-filtering-control/version_1_5_Julia
julia --project=. tools/profile_gpu.jl --both

# Python+JAX microbench (same surface, same metrics)
cd /home/ajay/Repos/python-smc2-filtering-control/version_1_5_Python_JAX
JAX_ENABLE_X64=True PYTHONPATH=.:.. \
  python tools/profile_gpu.py --both
\end{lstlisting}

For deeper Python+JAX profiling, pass \texttt{--trace} to dump a
\texttt{jax.profiler} trace into \texttt{./jax\_trace\_filter/}
and \texttt{./jax\_trace\_controller/} that TensorBoard's Profile
tab can render.

For deeper Julia profiling, run the bench under
\texttt{nsys profile --trace=cuda julia tools/bench\_smc\_full\_mpc\_fsa\_gpu.jl ...}
and inspect the resulting \texttt{.qdrep} in NVIDIA Nsight Systems.

\section{Repository basics: where things live}

\paragraph{Python+JAX.} Per-file layout under
\texttt{version\_1\_5\_Python\_JAX/}:

\begin{itemize}[itemsep=2pt]
  \item \texttt{models/fsa\_high\_res/\_dynamics.py} --- pure JAX
        drift, diffusion, EM step.
  \item \texttt{models/fsa\_high\_res/simulation.py} ---
        \texttt{DEFAULT\_PARAMS}, \texttt{INIT\_STATE},
        \texttt{params\_v15\_to\_v1}.
  \item \texttt{models/fsa\_high\_res/\_plant.py} --- pure
        \texttt{plant\_step} / \texttt{plant\_rollout}.
  \item \texttt{models/fsa\_high\_res/estimation.py} ---
        \texttt{HIGH\_RES\_FSA\_V15\_ESTIMATION}
        (\texttt{EstimationModel}); the Kalman-fused
        \texttt{\_propagate\_fn\_framework}.
  \item \texttt{models/fsa\_high\_res/control.py} --- RBF schedule
        + cost; \texttt{build\_control\_spec}.
  \item \texttt{tools/bench\_smc\_full\_mpc\_fsa\_v15.py} --- the
        closed-loop bench.
\end{itemize}

The \texttt{smc2fc} framework lives at \texttt{smc2fc/} (one level
up from any version dir). The relevant entry points consumed by
the bench are: \texttt{smc2fc.core.jax\_native\_smc.run\_smc\_window\_native},
\texttt{...run\_smc\_window\_bridge\_native},
\texttt{smc2fc.filtering.gk\_dpf\_v3\_lite.make\_gk\_dpf\_v3\_lite\_log\_density\_compileonce},
\texttt{smc2fc.control.tempered\_smc\_loop.run\_tempered\_smc\_loop\_native}.

\paragraph{Julia.} Per-file layout under
\texttt{version\_1\_5\_Julia/}:

\begin{itemize}[itemsep=2pt]
  \item \texttt{models/fsa\_high\_res/\_dynamics.jl} --- v1
        verbatim G0 SDE drift / diffusion / substepped EM.
  \item \texttt{models/fsa\_high\_res/simulation.jl} ---
        \texttt{DEFAULT\_PARAMS},
        \texttt{params\_v15\_to\_v1\_nt}, \texttt{sample\_obs\_bfa}.
        Uses \texttt{Match.jl} \texttt{@match} at every site that
        a future Lean4 port would also need to pattern-match.
  \item \texttt{models/fsa\_high\_res/\_plant.jl} --- pure
        \texttt{plant\_step} / \texttt{plant\_rollout}.
  \item \texttt{models/fsa\_high\_res/gpu\_pf.jl} ---
        \texttt{FSAGPUTarget}, \texttt{gpu\_log\_density},
        \texttt{gpu\_grads}, \texttt{parallel\_hmc\_one\_move},
        \texttt{parallel\_hmc\_one\_move!} (the new mutating
        variant).
  \item \texttt{models/fsa\_high\_res/gpu\_control.jl} ---
        \texttt{FSAv1ControlGPUTarget},
        \texttt{gpu\_cost\_log\_density\_batched}.
  \item \texttt{models/fsa\_high\_res/estimation.jl} ---
        \texttt{PARAM\_NAMES}, \texttt{PARAM\_PRIOR\_CONFIG}.
  \item \texttt{tools/bench\_smc\_full\_mpc\_fsa\_gpu.jl} --- the
        closed-loop bench, including \texttt{run\_outer\_smc}
        (the hand-rolled tempered SMC$^2$ loop) and
        \texttt{controller\_plan} (calls
        \texttt{SMC2FC.run\_tempered\_smc\_gpu}).
\end{itemize}

The Julia framework lives at \texttt{julia/SMC2FC/} (existing) and
\texttt{julia/SMC2FC\_functional/} (new, drop-in replacement).

\paragraph{Lean4 reference.} \texttt{version\_1\_5\_LEAN/}
provides a Lean4 implementation of the v1.5 model functions
exposed via a JSON-line CLI (\texttt{Main.lean}). The diff test
\texttt{version\_1\_5\_Julia/diff\_test/test\_lean\_diff\_v15.jl}
round-trips 121 randomised cases through the Lean4 binary
at 1e-6 single-step / 1e-4 integrated tolerance. This guarantees
the model math is identical across all three implementations.

\section{Closing notes}

The work in this session is committed and pushed on the
\texttt{Julia-port-verison-1} branch. Notable hashes:

\begin{itemize}[itemsep=2pt]
  \item \texttt{0c92907} --- CLI parity + GPU profilers (Phases A/A1/A1b)
  \item \texttt{ff3da57} --- per-stride telemetry + launcher + comparison script
  \item \texttt{155beec} --- T=7d pre-flight artefacts + initial findings
  \item \texttt{159f93e} --- Python A1 + A2 fix (shrinking horizon + mid-array splice)
  \item \texttt{3ff65c3} --- Python Kalman-fusion fix + \texttt{init\_state} cheat fixed
  \item \texttt{da5e9c9} --- findings reframing (drop misleading ``Python best of 3'')
  \item \texttt{c4a9e86} --- Julia perf: mutating HMC + optional Liu--West
  \item \texttt{ebbd9ed} --- v1.5 LEAN reference + 121-case round-trip diff test
  \item \texttt{a6903f9} --- \texttt{SMC2FC\_functional} library (parallel rewrite + 7-gate audit + GPU drop-in test)
  \item \texttt{da847d9} --- v1.5 Julia GPU filter: per-segment GPU-accumulator fix; \texttt{--ot-max-weight} CLI flag
  \item \texttt{099a14a} --- v1.5 Julia: default OT rescue OFF; heavy warnings on the slow path
  \item \texttt{4a5c007} --- \texttt{SMC2FC\_functional} migration on v1.5 GPU bench: bit-identical numerics
\end{itemize}

The corresponding plan archive is
\texttt{claude\_plans/FSA\_v1\_5\_three\_way\_closed\_loop\_comparison\_2026-05-08\_1423.md}
(also copied to \texttt{docs/PLAN\_ARCHIVE.md} in this directory).

\end{document}
